\section{Restricted M\"obius sums and their moments}\label{sec:arithmetic}

This section develops the arithmetic tools needed for the exceptional singular
value. All analytic cancellation comes from classical Mertens and smooth-number
estimates. The moment and support-constrained statements are consequences of
those inputs; no arithmetic priority claim is attached to their formulation.
Throughout, $n\to\infty$, and we write
\[
 L=\log n,\qquad s=s(n)=\sqrt{L\log L},\qquad
 b_p(n)=R(n/p,p),\qquad E_P(n)=\sum_{p\le n}b_p(n)^2.
\]
All sums indexed by $p$ or $q$ as primes below are explicitly so indicated when
there could be ambiguity. Arguments of counting functions may be real, with the
usual floor convention. Constants may depend on fixed compact intervals and on
fixed moment parameters, but never on a prime varying in the stated interval.

\subsection{Classical inputs and their ranges}

Let $\rho$ be the Dickman function, defined by $\rho(v)=1$ for $0\le v\le1$,
$\rho(v)=0$ for $v<0$, and
\[
 v\rho'(v)=-\rho(v-1)\qquad(v>1).
\]
We use the following classical estimates, recording the needed specialization
rather than assuming an asymptotic for $R$ itself. The ordinary Mertens bound is
\begin{equation}\label{arith:mertens}
 |M(z)|\le C z\exp\left\{-c\frac{(\log z)^{3/5}}
 {(\log\log z)^{1/5}}\right\}
 \le C z e^{-c(\log z)^{7/12}}\qquad(z\ge z_0).
\end{equation}
The second inequality follows after changing $c,C$, since
$(\log z)^{1/60}/(\log\log z)^{1/5}\to\infty$.
For a modern explicit statement implying the first inequality, including a
logarithmic prefactor absorbable into the exponential, see
Lee--Leong~\cite[Theorem~1.1, (12)]{leeleong}.

Here is the precise smooth-number input. Fix $0<c_-<c_+<\infty$ and $D>0$.
Suppose
\[
 c_-s\le t=\log y\le c_+s,\qquad
 |\log x-L|\le Ds,\qquad
 \log x-L^{9/10}\le\log z\le\log x.
\]
Uniformly over these choices, with $v=\log z/t$,
\begin{equation}\label{arith:smooth-input}
 \Psi(z,y)=z\left\{\rho(v)+(\gamma-1)\frac{\rho'(v)}t
 +O\left(\rho(v)\left(\frac{\log(v+1)}t\right)^2\right)\right\},
\end{equation}
where $\gamma$ is Euler's constant. We also use, uniformly as $v\to\infty$,
\begin{equation}\label{arith:dickman-input}
 -\rho'(v)\asymp\rho(v)\log(v+1),\qquad
 |\rho''(v)|\ll\rho(v)\log^2(v+1),\qquad
 \log\rho(v)=-(1+o(1))v\log v.
\end{equation}
Only arguments tending uniformly to infinity occur in the derivative estimates,
so the low-argument changes of differentiability are irrelevant.

For clarity, these are specializations of established smooth-number theory,
not additional hypotheses. McNew~\cite[(7), (13), (14), (17)]{mcnew}
records the Hildebrand first-order estimate, the de Bruijn--Saias expansion,
and the derivative estimates; see also Saias~\cite{saias} and the survey
of Hildebrand--Tenenbaum~\cite{hildebrandtenenbaum}.
The first-order statement alone would not suffice here.
In the notation of that expansion, take two terms (the derivative order
parameter $k=1$). The coefficient of $\rho'(v)/t$ is $\gamma-1$,
and the relative remainder after the two terms is of order
$(\log(v+1)/t)^2$ in the present range.
The source range contains
\[
 y>\exp\{(\log\log z)^{5/3+\varepsilon}\},\qquad
 y>(\log z)^{1+\varepsilon},\qquad y<z,
\]
with any fixed sufficiently small $\varepsilon>0$.
The extra boundary inequalities for $k=1$ in (13) are
\[
 \frac v2\ge\frac{\log t}{t},\qquad
 v-1\ge\frac{\log t}{t}.
\]
Indeed our substitutions give $\log z=L(1-o(1))$,
$t\asymp\sqrt{L\log L}$, and
$v\asymp\sqrt{L/\log L}$. All these inequalities therefore hold uniformly
with a diverging margin. The exponentially small relative error in the
Saias estimate is $o((\log(v+1)/t)^2)$ here. Since
$\rho'(v)/(t\rho(v))=O(\log(v+1)/t)=o(1)$, the multiplicative
remainder in the source yields exactly the additive remainder in
\eqref{arith:smooth-input}. In particular, no exclusion of integer $v$ is
needed in this range. The derivative and ratio estimates recorded there,
together with the defining Dickman equation, give
\eqref{arith:dickman-input}.

We also use the prime number theorem and Chebyshev's upper bound
$\pi(z)\ll z/\log z$. The series application additionally uses
Mertens' product bound
\[
 \prod_{p\le y}(1-1/p)^{-1}\ll\log(2y).
\]
The weaker estimate
\begin{equation}\label{arith:harmonic-product}
 \prod_{p\le y}(1-1/p)^{-1}
 =\exp\{O(\log\log(3y))\}
\end{equation}
already suffices for the moment proof. Partial summation of Chebyshev gives
\[
 \sum_{p\le y}p^{-1}=O(\log\log(3y)),
\]
and higher prime powers have a bounded total contribution to the logarithm
of the product.

\subsection{Two exact convolutions and a uniform upper bound}

Let $S_y(a)$ be the indicator of $P^+(a)\le y$, including $a=1$.
The coefficient identity
$(\mu*S_y)(d)=\mu(d)\one_{P^-(d)>y}$ follows prime by prime:
the local factor is $1$ at primes at most $y$ and $1-X$ at larger primes.
Summing its coefficients in either order gives
\begin{equation}\label{arith:convolutions}
 R(x,y)=\sum_{\substack{a\le x\\P^+(a)\le y}}M(x/a)
       =\sum_{k\le x}\mu(k)\Psi(x/k,y).
\end{equation}
These are finite identities and include the coefficient at $1$.
Replacing $M$ by $|M|$ in the first identity gives a positive majorant
which increases with $y$.

\begin{lemma}\label{arith:upper-local}
Fix $0<c_-<c_+<\infty$ and $D>0$. Uniformly for
$|\log x-L|\le Ds$ and $\log y=cs$ with $c_-\le c\le c_+$,
\[
 |R(x,y)|\le x\exp\left\{-\left(\frac1{2c}+o(1)\right)s\right\}.
\]
The same bound holds for the positive majorant just described.
\end{lemma}
\begin{proof}
Put $h=L^{9/10}$ and split the first sum in
\eqref{arith:convolutions} at $a=xe^{-h}$. By
\eqref{arith:mertens} with exponent $7/12$, the part below this point is
at most
\[
 Cxe^{-ch^{7/12}}\sum_{P^+(a)\le y}\frac1a
 \le x\exp\{-cL^{21/40}+O(\log L)\}=x e^{-\omega(s)}.
\]
Here and below $e^{-\omega(s)}$ means a quantity bounded by
$e^{-As}$ for every fixed $A>0$ once $n$ is sufficiently large.
For the remaining part use $|M(x/a)|\le x/a$ and Rankin's inequality.
With $u=\log x/t$, $t=\log y$, and $\alpha=\log u/t$, it is at most
\begin{equation}\label{arith:rankin}
 xe^{-\alpha(\log x-h)}
 \prod_{p\le y}(1-p^{-1+\alpha})^{-1}.
\end{equation}
Uniformly in this range, $u\to\infty$ and $0<\alpha\le1/4$ eventually.
Higher prime-power terms in the logarithm of the product are $O(1)$,
and Chebyshev partial summation bounds the prime terms by
\[
 O(1)+O(e^{\alpha t}/t)
 +O\left(\int_{\log2}^t\frac{e^{\alpha v}}v\,dv\right)
 =O(\log t+e^{\alpha t})=O(\log L+u)=o(s).
\]
To see the integral bound, split at $1/\alpha$. Below this point use a
bounded exponential; above it substitute $w=\alpha v$ and use
$\int_1^T e^w/w\,dw\ll e^T$. Finally $\alpha h=o(s)$ and
\[
 \alpha\log x=u\log u=\left(\frac1{2c}+o(1)\right)s
\]
uniformly on the fixed compact interval. Substitution in
\eqref{arith:rankin} completes the proof.
\end{proof}

\subsection{A relative asymptotic at the saddle scale}

\begin{lemma}\label{arith:bridge}
Under the ranges of Lemma~\ref{arith:upper-local}, put
$t=\log y$ and $u=\log x/t$. Then
\begin{equation}\label{arith:rough-asymptotic}
 R(x,y)=(1+o(1))\frac{x\rho'(u)}t
\end{equation}
uniformly. In particular it is negative for all sufficiently large $n$
throughout this range.
\end{lemma}
\begin{proof}
Set $a_0=7/12$, $K=\lfloor\exp((\log L)^2)\rfloor$,
$H=\exp(L^{9/10})$, and $\beta=\log(u+1)/t$.
Here $\beta\asymp\sqrt{\log L/L}$. For $0\le w\le L^{9/10}$,
all arguments between $u-w/t$ and $u$ are $(1-o(1))u$.
Integrating the logarithmic derivative bound in
\eqref{arith:dickman-input} gives
\begin{equation}\label{arith:dickman-shift}
 \rho(u-w/t)\le\rho(u)e^{C\beta w}.
\end{equation}

First truncate the second identity in \eqref{arith:convolutions} without
losing M\"obius cancellation. Interchanging its tail sums gives exactly
\begin{equation}\label{arith:tail-exact}
 T_K:=\sum_{K<k\le x}\mu(k)\Psi(x/k,y)
 =\sum_{\substack{a\le x/K\\P^+(a)\le y}}
       \{M(x/a)-M(K)\}.
\end{equation}
An endpoint with $x/a=K$ contributes zero, so this formula is unchanged by
floors. The contribution from $M(K)$ is bounded using
\eqref{arith:smooth-input} and \eqref{arith:dickman-shift} by
\[
 |M(K)|\Psi(x/K,y)
 \ll x\rho(u)\exp\{-c(\log K)^{a_0}+C\beta\log K\}.
\]
For the terms with $K<x/a<H$, split $x/a$ into dyadic intervals
$[q,2q)$ starting at $K$, cutting the last one at $H$ if necessary.
There are at most $\Psi(x/q,y)$ eligible integers $a$ in one interval.
The Mertens bound and the same two smooth estimates therefore bound its
contribution by
\[
 Cx\rho(u)\exp\{-c'(\log q)^{a_0}+C\beta\log q\}.
\]
Uniformly for $\log K\le v\le\log H$,
\[
 \beta v^{1-a_0}\le
 O\left(L^{-1/2+(9/10)(1-a_0)}\sqrt{\log L}\right)=o(1).
\]
Thus the positive term in the exponent is absorbed. Even counting all
$O(L)$ possible dyadic intervals, their sum is at most
$x\rho(u)\exp\{-c''(\log K)^{a_0}\}$ after reducing $c''$,
since $(\log K)^{a_0}\asymp(\log L)^{7/6}\gg\log L$.

For $x/a\ge H$, use \eqref{arith:harmonic-product} to obtain
\[
 \sum_{\substack{P^+(a)\le y\\x/a\ge H}}|M(x/a)|
 \ll x e^{-c(\log H)^{a_0}}
        \prod_{p\le y}(1-1/p)^{-1}
 \le x\exp\{-cL^{21/40}+O(\log L)\}.
\]
By \eqref{arith:dickman-input}, $|\log\rho(u)|=O(s)$; since
$21/40>1/2$, this last bound is $o(x\rho(u)\beta)$.
The previous bounds have the same property because
$(\log L)^{7/6}\gg|\log\beta|$. We have proved
\begin{equation}\label{arith:tail-small}
 T_K=o(x\rho(u)\beta)
\end{equation}
uniformly.

It remains to expand only the short initial sum $k\le K$.
Since $\beta\log K=o(1)$, Taylor's theorem and
\eqref{arith:dickman-input}--\eqref{arith:dickman-shift} give
\begin{align*}
 \rho(u-\log k/t)&=\rho(u)-\frac{\log k}{t}\rho'(u)
       +O\bigl(\rho(u)\beta^2(\log k)^2\bigr),\\
 \frac{\rho'(u-\log k/t)}t&=\frac{\rho'(u)}t
       +O\bigl(\rho(u)\beta^2\log k\bigr).
\end{align*}
Let
$A_0(K)=\sum_{k\le K}\mu(k)/k$ and
$A_1(K)=\sum_{k\le K}\mu(k)\log k/k$.
Applying \eqref{arith:smooth-input} and summing its errors absolutely yields
\begin{align}\label{arith:short-expansion}
 \sum_{k\le K}\mu(k)\Psi(x/k,y)
 &=x\rho(u)A_0(K)-\frac{x\rho'(u)}t A_1(K)\notag\\
 &\quad +(\gamma-1)\frac{x\rho'(u)}t A_0(K)
 +O\bigl(x\rho(u)\beta^2(1+\log K)^3\bigr).
\end{align}
The cubic factor bounds $\sum_{k\le K}(1+\log k)^2/k$.

For completeness, the two M\"obius moments here satisfy
\begin{equation}\label{arith:mobius-moments}
 A_0(K)=O(e^{-c_1(\log K)^{a_0}}),\qquad
 A_1(K)=-1+O(e^{-c_1(\log K)^{a_0}}).
\end{equation}
Indeed partial summation with \eqref{arith:mertens} proves convergence
of both infinite series and bounds their tails by a polynomial in
$\log K$ times $e^{-c(\log K)^{a_0}}$.
The same integrable bounds justify letting $\sigma\downarrow1$ in
$\sum\mu(k)k^{-\sigma}=1/\zeta(\sigma)$ and its derivative.
The simple pole of $\zeta$ at $1$ then gives values $0$ and $-1$,
respectively. Absorb the polynomial into a smaller $c_1$ to obtain
\eqref{arith:mobius-moments}.
Now $\beta(1+\log K)^3=o(1)$ and
$|\rho'(u)|/t\asymp\rho(u)\beta$.
Equations \eqref{arith:tail-small}--\eqref{arith:mobius-moments}
leave precisely $x\rho'(u)/t$ as the main term, proving the lemma.
\end{proof}

The two-term input is essential to this argument. A relative error of order
$\beta$ in the first-order smooth estimate would be as large as the surviving
term after $A_0(\infty)=0$ cancels the density term.

\subsection{Moments, the maximum, and prime concentration}

\begin{theorem}\label{arith:moments}
For every fixed real $q>1$,
\begin{equation}\label{arith:moment-law}
 \sum_{p\le n}|b_p(n)|^q
 =n^q\exp\{-(\sqrt{2q(q-1)}+o(1))s\}.
\end{equation}
Moreover
\begin{equation}\label{arith:maximum}
 \max_{p\le n}|b_p(n)|
 =n\exp\{-(\sqrt2+o(1))s\}.
\end{equation}
For every fixed $0<\delta<1$, there is $\eta_\delta>0$ such that
\begin{equation}\label{arith:concentration}
 \sum_{\substack{p\le n\\|\log p/s-1|>\delta}}b_p(n)^2
 \le n^2e^{-(2+\eta_\delta)s}
\end{equation}
for all sufficiently large $n$.
\end{theorem}
\begin{proof}
For a fixed compact interval $0<c_-\le\log p/s\le c_+$,
Lemma~\ref{arith:upper-local} applied to $x=n/p$, $y=p$ gives
\[
 |b_p(n)|\le\frac np
 \exp\left\{-\left(\frac{s}{2\log p}+o(1)\right)s\right\}.
\]
In a bin $e^{as}<p\le e^{bs}$ within this interval, it follows that
\begin{equation}\label{arith:bin}
 \sum_{e^{as}<p\le e^{bs}}|b_p(n)|^q
 \ll_q n^q\exp\left\{-\left((q-1)a+\frac q{2b}-o(1)\right)s\right\},
\end{equation}
since $\sum_{m>e^{as}}m^{-q}\ll_q e^{-(q-1)as}$.
Cover the compact interval by finitely many bins of mesh at most $\varepsilon$.
As the mesh decreases their costs approach
\[
 f_q(c)=(q-1)c+\frac q{2c},\qquad
 \min_{c>0}f_q(c)=\sqrt{2q(q-1)},\qquad
 c_q=\sqrt{\frac q{2(q-1)}}.
\]

The two noncompact prime ranges require separate bounds. For $p<e^{c_-s}$,
enlarge the smoothness parameter in the positive majorant of
\eqref{arith:convolutions} to $y_0=e^{c_-s}$, keeping $x=n/p$.
Lemma~\ref{arith:upper-local} gives a total at most
\[
 n^q\exp\{-(q/(2c_-)-o(1))s\}\sum_p p^{-q}.
\]
Choose $c_-$ sufficiently small. For $p>e^{c_+s}$ the bound
$|b_p(n)|\le n/p$ gives total
$O_q(n^qe^{-(q-1)c_+s})$; choose $c_+$ sufficiently large.
For each desired error first fix these endpoints and a sufficiently fine
mesh, and only then let $n\to\infty$.
This proves the upper bound in \eqref{arith:moment-law}.
The same argument pointwise minimizes $c+1/(2c)$ to $\sqrt2$.
Its small-prime and large-prime costs are $1/(2c_-)$ and $c_+$,
respectively, proving the upper bound in \eqref{arith:maximum}
without taking a nonuniform limit $q\to\infty$.

For the lower bounds fix any $c>0$ and take primes in
$[P,2P]$, $P=e^{cs}$. Lemma~\ref{arith:bridge} and
\eqref{arith:dickman-input} give uniformly in this block
\begin{equation}\label{arith:block-law}
 |b_p(n)|=n\exp\{-(c+1/(2c)+o(1))s\}.
\end{equation}
Indeed $u=\log(n/p)/\log p$ satisfies
$u\log u=(1/(2c)+o(1))s$, and logarithms of the derivative and
$1/\log p$ prefactors are $o(s)$.
The prime number theorem supplies
$\pi(2P)-\pi(P)=\exp\{(c+o(1))s\}$ primes.
Taking $c=c_q$ yields the matching moment lower bound, while taking
$c=1/\sqrt2$ and any prime in the block yields the maximum lower bound.

Finally set $q=2$ in the bin argument. The cost is $c+1/c$,
whose unique minimum is $2$ at $c=1$.
Outside $(1-\delta,1+\delta)$ its infimum exceeds $2$.
Choose the remote endpoints so that their costs also exceed $2$,
and choose a mesh small enough to preserve half the positive gap.
The same finite-bin estimates, absorbing their number into the gap,
prove \eqref{arith:concentration}.
\end{proof}

No assertion here is uniform in a varying $q$. In particular,
\eqref{arith:moment-law} at $q=2$ gives
$E_P(n)=n^2e^{-(2+o(1))s}$, and \eqref{arith:concentration}
then gives concentration relative to the total prime energy.

\subsection{Composite multipliers and the inverse-vector norm}

Recall the full energy and the inverse-vector norm from the finite identities:
\[
 E(n)=\sum_{\substack{2\le j\le n\\\mu(j)^2=1}}
       R(n/j,P^+(j))^2,\qquad W_n=\norm{w_n}_2,
\]
\begin{equation}\label{arith:norm-recalled}
 W_n^2=E(n)+(M(n)-1)^2+n-Q(n).
\end{equation}
Here $Q(n)=\sum_{j\le n}\mu(j)^2$ is a count, not its square root.

\begin{theorem}\label{arith:energy}
The prime and full energies and the inverse-vector norm satisfy
\begin{gather}
 E_P(n)=n^2e^{-(2+o(1))s},\qquad
 E(n)=n^2e^{-(2+o(1))s},\qquad W_n=ne^{-(1+o(1))s},
 \label{arith:energy-scale}\\
 \frac{E(n)}{E_P(n)}\longrightarrow\frac{15}{\pi^2},\qquad
 \frac{W_n^2}{E_P(n)}\longrightarrow\frac{15}{\pi^2}.
 \label{arith:energy-ratio}
\end{gather}
These conclusions do not give a multiplicative equivalent for $E_P(n)$ itself.
\end{theorem}
\begin{proof}
The prime-energy assertion has just been proved. Every squarefree $j\ge2$
has a unique factorization $j=ap$ with $p=P^+(j)>P^+(a)$ and $a$
squarefree, including $a=1$. Consequently
\begin{equation}\label{arith:multiplier-decomposition}
 E(n)=\sum_{a\ge1}\mu(a)^2
       \sum_{\substack{p>P^+(a)\\ap\le n}}R(n/(ap),p)^2.
\end{equation}
Put $A=e^{3s}$ and $\mathcal C_n=\{p:e^{s/2}\le p\le e^{2s}\}$.
The contribution of $a>A$, by $|R(x,p)|\le x$ for $x\ge1$, is at most
\[
 n^2\sum_{a>A}a^{-2}\sum_p p^{-2}=O(n^2/A)=o(E_P(n)).
\]

For $a\le A$ set $N_a=\lfloor n/a\rfloor$. The exact floor identity is
\begin{equation}\label{arith:floor}
 R(n/(ap),p)=R(N_a/p,p):
 \quad pd\le n/a\ \Longleftrightarrow\ pd\le N_a
 \quad\hbox{for integers }p,d.
\end{equation}
Uniformly over these multipliers,
\[
 N_a=(n/a)(1+O(A/n)),\qquad
 \log N_a=L-\log a+O(A/n),\qquad s(N_a)/s=1+o(1),
\]
and $\min_{a\le A}N_a\to\infty$.
For all sufficiently large $n$, the complement of $\mathcal C_n$ lies
outside the band $|\log p/s(N_a)-1|\le1/3$, uniformly in $a\le A$.
Apply \eqref{arith:concentration} at the integer $N_a$ and drop the
restriction $p>P^+(a)$. For one fixed $\eta>0$, the total contribution
from $a\le A$, $p\notin\mathcal C_n$ is bounded by
\[
 \sum_{a\le A}N_a^2e^{-(2+\eta)s(N_a)}
 \le n^2e^{-(2+\eta/2)s}\sum_{a\ge1}a^{-2}=o(E_P(n)).
\]
This invocation is uniform because the same theorem holds at every
sufficiently large integer, and all the $N_a$ tend to infinity together.

It remains to compare central terms multiplicatively. For $p\in\mathcal C_n$
put $t=\log p$, $u=\log(n/p)/t$, and $v=u-\log a/t$.
Lemma~\ref{arith:bridge} applies uniformly both at $(n/p,p)$ and at
$(n/(ap),p)$ because their first logarithms equal $L+O(s)$.
Equivalently one may use \eqref{arith:floor}: replacing $N_a$ by $n/a$
changes the Dickman argument by $O(A/(nt))$, which is negligible under
the logarithmic derivative bound below. Thus
\begin{equation}\label{arith:multiplier-ratio}
 \frac{R(n/(ap),p)}{R(n/p,p)}
 =(1+o(1))\frac1a\frac{\rho'(v)}{\rho'(u)}
 \qquad(a\le A,\ p\in\mathcal C_n)
\end{equation}
with a common relative error. The denominator is nonzero by
Lemma~\ref{arith:bridge}. All arguments from $v$ to $u$ tend uniformly
to infinity, and $0\le u-v\le6$.
Equation \eqref{arith:dickman-input} gives
\[
 \left|\frac{d}{dz}\log(-\rho'(z))\right|
 \ll\log(z+1),\qquad
 \left|\frac{\rho'(v)}{\rho'(u)}\right|
 \le\exp\left\{C\frac{\log L}{s}\log a\right\}.
\]
As $\log L/s\to0$, for any fixed $0<\varepsilon<1$ we obtain
\begin{equation}\label{arith:domination}
 R(n/(ap),p)^2\le C a^{-2+\varepsilon}R(n/p,p)^2
\end{equation}
uniformly in the truncated region for all sufficiently large $n$.
For each fixed $a$, \eqref{arith:multiplier-ratio} instead gives the
uniform equivalent $a^{-2}$ after squaring; also $p>P^+(a)$ throughout
$\mathcal C_n$ eventually.

Define the normalized central contribution explicitly by
\[
 \mathcal E_{a,n}=\mathbf1_{a\le A}\frac{\mu(a)^2}{E_P(n)}
 \sum_{\substack{p\in\mathcal C_n\\p>P^+(a),\ ap\le n}}
 R(n/(ap),p)^2.
\]
It is bounded by
$C\mu(a)^2a^{-2+\varepsilon}$, a summable sequence independent of $n$.
For each fixed $a$ it tends to $\mu(a)^2/a^2$, since central prime energy
is $(1+o(1))E_P(n)$. Dominated convergence, followed by the two discarded
tail estimates, proves
\[
 \frac{E(n)}{E_P(n)}\longrightarrow
 \sum_{a\ge1}\frac{\mu(a)^2}{a^2}
 =\prod_p(1+p^{-2})=\frac{\zeta(2)}{\zeta(4)}=\frac{15}{\pi^2}.
\]
Finally \eqref{arith:mertens} implies
$(M(n)-1)^2=o(E_P(n))$, since $L^{7/12}/s\to\infty$,
and $n-Q(n)\le n=o(E_P(n))$. The exact norm identity
\eqref{arith:norm-recalled} proves the remaining assertions.
\end{proof}

\subsection{Sharp support constraints and exceptional cutoffs}

For fixed $\kappa>0$ define
\[
 F_\kappa(n)=\sup_{\substack{|a_p|\le1\\
 \#\{p\le n:a_p\ne0\}\le\lfloor e^{\kappa s}\rfloor}}
 \left|\sum_{2\le m\le n}\mu(m)a_{P^-(m)}\right|.
\]
Weights may be complex and may depend on $n$. Grouping a squarefree
integer $m$ by its least prime $p$ gives the exact identity
\begin{equation}\label{arith:least-prime}
 \sum_{2\le m\le n}\mu(m)a_{P^-(m)}
 =-\sum_{p\le n}a_p R(n/p,p)=-\sum_{p\le n}a_p b_p(n),
\end{equation}
since $m=pd$ has $P^-(d)>p$ and $\mu(pd)=-\mu(d)$.
Alladi~\cite{alladi1982} treats the unrestricted bounded-weight problem;
here the support ceiling is part of the extremal question.

\begin{theorem}\label{arith:sparse}
For every fixed $\kappa>0$,
\begin{equation}\label{arith:sparse-law}
 F_\kappa(n)=n\exp\{-(g(\kappa)+o(1))s\},\qquad
 g(\kappa)=
 \begin{cases}
 \sqrt2-\kappa,&0<\kappa\le1/\sqrt2,\\
 1/(2\kappa),&\kappa\ge1/\sqrt2.
 \end{cases}
\end{equation}
\end{theorem}
\begin{proof}
If $K\le e^{\kappa s}$ weights are nonzero, H\"older's inequality and
\eqref{arith:moment-law} imply, for each fixed $q>1$,
\[
 \left|\sum_pa_pb_p(n)\right|
 \le K^{1-1/q}\left(\sum_p|b_p(n)|^q\right)^{1/q}
 \le n\exp\{-(\sqrt{2\theta}-\kappa\theta-o(1))s\},
 \qquad\theta=1-1/q.
\]
For $\kappa>1/\sqrt2$, choose the fixed
$\theta=1/(2\kappa^2)\in(0,1)$ to obtain saving $1/(2\kappa)$.
For $\kappa\le1/\sqrt2$, use directly $K\max_p|b_p(n)|$
and \eqref{arith:maximum} to obtain saving $\sqrt2-\kappa$.
This covers the common boundary without varying $q$.

For the lower bound when $\kappa<1/\sqrt2$, take
$\lfloor e^{\kappa s}\rfloor$ primes from
$[e^{s/\sqrt2},2e^{s/\sqrt2}]$. The prime number theorem guarantees
at least this many primes for all sufficiently large $n$.
Equation \eqref{arith:block-law} gives each selected coefficient magnitude
$n e^{-(\sqrt2+o(1))s}$. Choose $a_p$ to align their contributions
in \eqref{arith:least-prime}; their sum gives the required lower bound.
For $\kappa\ge1/\sqrt2$, take every prime in
$[e^{\kappa s},2e^{\kappa s}]$. Their number is asymptotic to
$e^{\kappa s}/(\kappa s)$, hence satisfies the support ceiling, and their
total magnitude by \eqref{arith:block-law} is
$n e^{-(1/(2\kappa)+o(1))s}$. This also proves the boundary case.
In fact the coefficients in either block are eventually negative by
Lemma~\ref{arith:bridge}, so indicator weights suffice for these lower
constructions. The supremum assertion does not claim sharpness for every
prescribed support.
\end{proof}

\begin{corollary}\label{arith:exceptional}
Let $N_\lambda(n)=\#\{p\le n:|b_p(n)|\ge ne^{-\lambda s}\}$.
For fixed $\lambda>\sqrt2$,
\[
 N_\lambda(n)=\exp\{(h(\lambda)+o(1))s\},\qquad
 h(\lambda)=\frac{\lambda+\sqrt{\lambda^2-2}}2.
\]
For fixed $\lambda<\sqrt2$ the count is eventually zero.
At $\lambda=\sqrt2$ we obtain only the upper bound
$N_{\sqrt2}(n)\le\exp\{(1/\sqrt2+o(1))s\}$.
\end{corollary}
\begin{proof}
Markov's inequality and \eqref{arith:moment-law} give, for fixed $q>1$,
\[
 N_\lambda(n)\le
 \exp\{(q\lambda-\sqrt{2q(q-1)}+o(1))s\}.
\]
For $\lambda>\sqrt2$ the minimum is attained at
$q=(1+\lambda/\sqrt{\lambda^2-2})/2$, and its value is $h(\lambda)$.
For the reverse bound choose a fixed $c$ strictly between the two roots
of $c+1/(2c)=\lambda$ and arbitrarily close to the larger root
$h(\lambda)$. Every prime in $[e^{cs},2e^{cs}]$ meets the threshold
eventually by \eqref{arith:block-law}, so the count is at least
$e^{(c+o(1))s}$. Let $c$ increase to $h(\lambda)$ after taking the
limit in $n$. The assertion for $\lambda<\sqrt2$ follows from
\eqref{arith:maximum}. At equality the displayed Markov bound holds
for each fixed $q$, and its cost decreases to $1/\sqrt2$ as
$q\to\infty$. Taking the infimum after the limsup proves the stated
upper bound; no uniform moment estimate or matching lower is asserted.
\end{proof}

\subsection{A weighted harmonic series on a thin prime support}

\begin{theorem}\label{arith:thin-series}
Fix $\kappa>0$ and a fixed set $\mathcal S$ of primes such that
\[
 N_{\mathcal S}(x):=\#(\mathcal S\cap[2,x])\le e^{\kappa s(x)}
\]
for every sufficiently large $x$. For fixed complex weights $|a_p|\le1$
supported on $\mathcal S$, the series in increasing order of $m$ satisfies
\begin{equation}\label{arith:series-value}
 \sum_{m\ge2}\frac{\mu(m)a_{P^-(m)}}m=0,
 \qquad
 \left|\sum_{m>X}\frac{\mu(m)a_{P^-(m)}}m\right|
 \le\exp\{-(g(\kappa)-o(1))s(X)\}.
\end{equation}
If $\log(N_{\mathcal S}(x)+1)=o(s(x))$, the tail bound holds with
$g(\kappa)$ replaced by $\sqrt2$.
\end{theorem}
\begin{proof}
Write $c_m=\mu(m)a_{P^-(m)}$ for $m\ge2$ and
$T(x)=\sum_{2\le m\le x}c_m$.
At every sufficiently large integer endpoint the support hypothesis and
Theorem~\ref{arith:sparse} apply, and passage to real $x$ by flooring
is harmless. For each fixed $0<\varepsilon<g(\kappa)/2$,
\[
 |T(x)|\le x e^{-(g(\kappa)-\varepsilon)s(x)}
\]
eventually. For every fixed $c>0$, substitution $v=\log t$ and then
$z=\sqrt{v\log v}$ shows
\begin{align*}
 \int_X^\infty e^{-cs(t)}\frac{dt}{t}
 &=\int_{\log X}^\infty e^{-c\sqrt{v\log v}}\,dv\\
 &\le\int_{s(X)}^\infty 2z e^{-cz}\,dz
 =O_c((s(X)+1)e^{-cs(X)}),
\end{align*}
using $dv/dz=2z/(\log v+1)\le2z$ eventually.
Partial summation now proves convergence and bounds the tail by
$|T(X)|/X+\int_X^\infty|T(t)|t^{-2}\,dt$.
Absorbing the polynomial factor into an arbitrarily small exponential
loss proves the tail bound in \eqref{arith:series-value}.

Convergence alone does not identify the value. For real $\sigma>1$,
absolute convergence allows grouping by the least prime and Euler products:
\begin{equation}\label{arith:series-euler}
 \sum_{m\ge2}\frac{c_m}{m^\sigma}
 =-\frac1{\zeta(\sigma)}
   \sum_{p\in\mathcal S}\frac{a_p}{p^\sigma}
      \prod_{\substack{q\le p\\q\ \mathrm{prime}}}(1-q^{-\sigma})^{-1}.
\end{equation}
The support assumption implies $N_{\mathcal S}(x)\le x^{1/2}$
eventually. Partial summation of this estimate gives
$\sum_{p\in\mathcal S}\log(2p)/p<\infty$: the integral of
$x^{1/2}\log(2x)/x^2$ converges, and its boundary term tends to zero.
For $\sigma\ge1$, Mertens' product bound dominates the absolute value
of the $p$-summand on the right of \eqref{arith:series-euler} by
$C\log(2p)/p$. Thus the prime sum remains bounded as
$\sigma\downarrow1$, whereas $1/\zeta(\sigma)\to0$.
The Dirichlet-series limit is zero. Finally partial summation writes the
left side as
$\sigma\int_1^\infty T(t)t^{-\sigma-1}\,dt$.
The bound already proved makes $|T(t)|/t^2$ integrable, so dominated
convergence for $1<\sigma\le2$ identifies its limit with the convergent
series at $\sigma=1$. Its value is therefore zero.

Under the smaller-support assumption,
\eqref{arith:least-prime} and \eqref{arith:maximum} directly give
$|T(x)|\le x e^{-(\sqrt2-o(1))s(x)}$.
The support still satisfies the preceding polynomial counting bound,
and the identical integral and Dirichlet-series arguments yield the
last assertion.
\end{proof}

The theorem concerns the naturally ordered series, not absolute convergence
of its terms, and supplies no explicit numerical onset. The unweighted
zero-value assertion for a prime set of density zero is already contained in
the density formula of Kural--McDonald--Sah~\cite[Theorem~1.1]{kms}.
The formulation here records a quantitative tail for fixed thin supports
and bounded complex weights; it does not assert that this rate has no prior
occurrence in the literature.
